\section*{第 \ref{ch:g1:materials-around-us} 章——物体是由什么材料制成的}

\begin{solution}{exo:g1:materials-around-us:1}
椅子是一个物体。木头是一种材料：椅子可以用它做成。
\end{solution}

\begin{solution}{exo:g1:materials-around-us:2}
窗玻璃：玻璃。钉子：金属。鹅卵石：石头。书：纸。袜子：
织物（羊毛或棉花）。
\end{solution}

\begin{solution}{exo:g1:materials-around-us:3}
木头：铅笔、门（或椅子、桌子）。塑料：碗、水桶
（或玩具、瓶子）。
\end{solution}

\begin{solution}{exo:g1:materials-around-us:4}
金属叉子：它是唯一不是用木头或塑料做成的。（两个用木头做成，
一个用塑料做成。）
\end{solution}

\begin{solution}{exo:g1:materials-around-us:5}
玻璃：它是透明的。
\end{solution}

\begin{solution}{exo:g1:materials-around-us:6}
钢钉。磁铁能吸引铁和钢，不能吸引玻璃或木头。
\end{solution}

\begin{solution}{exo:g1:materials-around-us:7}
两个物体漂浮（木块和塑料瓶盖）；三个下沉
（鹅卵石、弹珠和钉子）。
\end{solution}

\begin{solution}{exo:g1:materials-around-us:8}
刀片是金属的，手柄是塑料的。剪刀由两种
材料制成。
\end{solution}

\begin{solution}{exo:g1:materials-around-us:9}
在蓝圈里的：它是金属做的。它不在橙圈里，
因为磁铁吸不住它：它是铝做的，不是铁
或钢做的。
\end{solution}

\begin{solution}{exo:g1:materials-around-us:10}
窗户必须让光透进来，我们必须能透过它看见外面。玻璃
是透明的；木头不是。
\end{solution}

\begin{solution}{exo:g1:materials-around-us:11}
“是”：钉子、回形针和钥匙，3 个物体。“否”：
鹅卵石、勺子、罐子、弹珠和袜子，5 个物体。
$3 + 5 = 8$：每个物体都各有其位。
\end{solution}
